% Notes of the Montreal seminar, 1970 (cote 129, typed part): archive cote 129, site page 44 (PDF page 45).
% typescript. Draft compared: none.
% AI-transcribed by Claude Opus 5.5 (claude-opus-5-5, adaptive thinking, effort high), Message Batches API, batch mtl129_full (msgbatch_01AGnDKHwy6ATCM7ehUa6263), 2026-11-13; not yet checked by eye.
\scanpage{129/44}
\mspage{\hand{II.16}}

s'identifie à $F_{D^*(G)}$ car $u \circ \mathbb{F}_{G/k} = \mathbb{F}_{\underset{\rightarrow}{W}} \circ u$ et $\mathbb{F}_{\underset{\rightarrow}{W}}$ n'est autre que
$F_{\underset{\rightarrow}{W}}$ (I, 5.3).

On fait un raisonnement analogue pour établir l'assertion
relative à $D^*(\mathbb{V}_G)$ .

Si $G$ est unipotent, il résulte de ce que $\mathbb{V}_G$ est nilpotent
(2.5.2) que $V_{D^*(G)}$ opère de façon nilpotente. On peut alors montrer,
plus généralement, que le foncteur $D^*$ définit une équivalence de ca-%
tégories entre la catégorie des p-groupes affines unipotents sur $k$ avec
la catégorie des \hand{$\mathcal{D}$}-modules sur lesquels $V$ est localement nilpotent. On
a en particulier :

\uline{Corollaire 4.3.3} On a des équivalences de catégories :

\diagram{\[
\left(\begin{array}{l} \text{p-groupes finis} \\ \text{unipotents sur } k \end{array}\right)^0 , \xrightarrow[\hand{\approx}]{D^*} \left(\begin{array}{l} \text{sous-catégorie pleine de } \underline{\mathrm{Mod}}_{\hand{\mathcal{D}},W\,\hand{\ell.f}} \\ \text{formée des objets } M \text{ où } V \text{ est nilpotent} \end{array}\right)
\]}

\diagram{\[
\left(\begin{array}{l} \text{p-groupes finis} \\ \text{étales sur } k \end{array}\right)^0 \xrightarrow[\hand{\approx}]{D^*} \left(\begin{array}{l} \text{sous-catégorie pleine de } \underline{\mathrm{Mod}}_{\hand{\mathcal{D}},W\,\hand{\ell.f.}} \\ \text{formée des objets } M \text{ où } F \text{ est un isomorphisme} \end{array}\right.
\]}

\diagram{\[
\left(\begin{array}{l} \text{p-groupes finis} \\ \text{bi-infinitésimaux} \end{array}\right)^0 \xrightarrow[\hand{\approx}]{D^*} \left(\begin{array}{l} \text{sous-catégorie pleine de } \underline{\mathrm{Mod}}_{\hand{\mathcal{D}},W\,\hand{\ell.f.}} \\ \text{formée des objets où } F \text{ et } V \text{ sont nilpotents} \end{array}\right)
\]}

Cela résulte de la description de $F$ et $V$ (4.3.2) et de
la caractérisation de chaque type de groupes à l'aide des morphismes $\mathbb{F}$
et $\mathbb{V}$ . (2.5).

\margin{44}
